%%%%%%%%%%%%%%%%%%%%%%%%%%%%%%%%%%%%%%%%%%%%%%%%%%%%%%%%%%%%%%%%%%%%%%
%  月度总结 · 阶段二（研究推进）  ——  空白骨架
%
%  怎么用
%    1. 把文件另存为  2026-10_阶段二_你的名字_方向.tex
%    2. 填下面「只改这里」的六项信息
%    3. 在每个章节的 { } 或环境里填入你的内容
%    4. 编译：xelatex 文件名.tex   （跑两遍）
%       Overleaf：Menu -> Compiler -> XeLaTeX
%
%  阶段二与阶段一的区别
%    · 主线是「研究推进」而不是「学习过程」
%    · 「关键结果」必须有量化数据
%    · 「与预期的差异」和「尚不能确定的」不能为空
%      —— 全都符合预期等于没有深入思考
%%%%%%%%%%%%%%%%%%%%%%%%%%%%%%%%%%%%%%%%%%%%%%%%%%%%%%%%%%%%%%%%%%%%%%
\documentclass[stage2]{ydzjreport}

% ===================== 只改这里：你的信息 =====================
\stagetwosetup
  {阶段二 · 进阶阅读}%          1 报告阶段
  {你的姓名}%                   2 姓名
  {你的方向}%                   3 方向（2-6 个汉字）
  {2026-10-01 至 2026-10-31}%   4 报告周期
  {关键词一；关键词二；关键词三}%   5 关键词
  {一句话说清本月推进了什么。}%     6 一句话摘要
% ============================================================

\begin{document}
\makestagetwohead

% ==================================================================
\section{本月研究任务}

\subsection{任务目标}

% 月初定的目标，2--4 条。
\begin{enumerate}
  \item 
  \item 
\end{enumerate}

\subsection{完成情况}

% 状态列只写三种：已完成 / 部分完成 / 未开始。
\begin{table}[htbp]
\centering
\caption{本月任务完成情况}
\begin{tabularx}{\textwidth}{@{}Y c Y@{}}
\toprule
\rowcolor{ydzjlight} 任务 & 状态 & 说明 \\
\midrule
 & 已完成 &  \\
 & 部分完成 &  \\
 & 未开始 &  \\
\bottomrule
\end{tabularx}
\end{table}


% ==================================================================
\section{进阶阅读与理论支撑}

\subsection{本阶段文献清单}

% 与阶段一不同：这里每篇文献必须对应「解决了我哪个问题」。
\begin{table}[htbp]
\centering
\caption{本阶段文献}
\begin{tabularx}{\textwidth}{@{}c Y Y@{}}
\toprule
\rowcolor{ydzjlight} 序号 & 文献 & 解决了我哪个问题 \\
\midrule
1 &  &  \\
2 &  &  \\
\bottomrule
\end{tabularx}
\end{table}

\subsection{对本方向的启示}

% 这些文献怎么改变了我的做法。150--300 字。
% 关键：要写出「因为读了 X，所以我决定改成 Y」。


% ==================================================================
\section{研究进展}

\subsection{已完成的工作}

% 具体做了什么，写到别人可以照着做的程度。
\begin{enumerate}
  \item 
  \item 
\end{enumerate}

\subsection{关键结果}

% ★ 必须有量化数据：表格或图。不能只有文字描述。
\begin{table}[htbp]
\centering
\caption{关键结果}
\begin{tabularx}{\textwidth}{@{}Y c Y@{}}
\toprule
\rowcolor{ydzjlight} 指标 & 数值 & 说明 \\
\midrule
 &  &  \\
 &  &  \\
\bottomrule
\end{tabularx}
\end{table}

\subsection{与预期的差异}

% ★ 本节不可为空。即使结果符合预期，也要写清「为什么符合」。
%   有偏差就分析原因（把相关曲线/数据叠在一起看，找主因）。


% ==================================================================
\section{问题与对策}

\subsection{遇到的主要问题}

% 至少 2 条，写清卡在哪。
\begin{enumerate}
  \item 
  \item 
\end{enumerate}

\subsection{已尝试的方案}

% 试过什么、为什么不行。与上面的问题一一对应。
\begin{enumerate}
  \item 
  \item 
\end{enumerate}

\subsection{下一步对策}

% 打算怎么办。与上面的问题一一对应。
\begin{enumerate}
  \item 
  \item 
\end{enumerate}


% ==================================================================
\section{阶段性结论}

\subsection{可以确定的}

% 有证据支撑的结论，至少 2 条。
\begin{enumerate}
  \item 
  \item 
\end{enumerate}

\subsection{尚不能确定的}

% ★ 本节不可为空，至少 1 条。写清「缺什么证据才能确定」。
\begin{enumerate}
  \item 
\end{enumerate}


% ==================================================================
\section{下月计划}

\begin{table}[htbp]
\centering
\caption{下月工作计划}
\begin{tabularx}{\textwidth}{@{}Y Y c@{}}
\toprule
\rowcolor{ydzjlight} 事项 & 预期产出 & 时间 \\
\midrule
 &  & 第 1 周 \\
 &  & 第 2--3 周 \\
 &  & 第 4 周 \\
\bottomrule
\end{tabularx}
\end{table}


% ==================================================================
\begin{thebibliography}{9}
\bibitem{ref1} 
\bibitem{ref2} 
\end{thebibliography}

\end{document}
